\section{Native weighted normal components and essential-disc counts}
\label{sec:weighted-components}

The weighted component operations of report 53 are now integrated into the
maintained package. They answer a question unavailable from the earlier
aggregate normal-surface statistics: how many components of this supplied
surface are compressing discs? Binary normal coordinates stay compressed.
The full component census and the reduced count-only query have independent
source-bound certificate replay. They require the existing validated finite,
connected, orientable triangulation with one torus boundary. A positive answer
concerns that manifold and vector; relating the manifold to an input knot
exterior remains a separate provenance obligation.

\subsection{Transport additive weights along an orbit trace}

Assign a vector $w(x)\in\mathbb Z^d$ to each point of an interval-pairing
universe $[0,N)$. The input is an additive list of half-open intervals with
constant vectors, possibly overlapping and with signed entries. Unmarked
points have weight zero. Endpoint differences combine the intervals into
canonical constant runs. The desired compressed histogram records
\[
 H(v)=\#\left\{O\text{ an orbit}:\sum_{x\in O}w(x)=v\right\}.
\]
Equal vectors share one row with binary multiplicity. Negative Euler weights
are ordinary signed integers; the multiplicities remain nonnegative counts.

The producer first obtains a complete trace from the maintained AHT engine.
It then transports weights through its local moves. Identity deletion,
trimming, periodic merging and transmission alter the description of the
same equivalence relation and require no weight change. Truncation removes
points after transferring their weights to surviving representatives.
Contraction emits the weights of static singleton representatives and closes
their gaps. At the end, every original orbit has one emitted representative
carrying its complete sum. The total count and every coordinate's total
weighted mass are checked against the original source.

For a translation truncation with new endpoint $c$ and period $p$, set
$b=c-p$. A removed point $x$ maps to
\[
 b+((x-b)\bmod p)\in[b,c).
\]
For a constant-weight deleted interval of length $\ell=qp+r$, add $q$
copies of its vector to the whole period and one copy to the cyclic remainder,
which has at most two ordinary interval pieces. Keep the preexisting prefix
weights as well. All arithmetic uses the binary quotient, remainder and
endpoints; no loop has $q$ iterations. For a reflection $x\mapsto a+N-1-x$,
a removed half-open interval $[u,v)$ contributes its weight to
$[a+N-v,a+N-u)$. This reverses the interval, with the endpoint convention
accounted for exactly.

The independent weighted verifier first checks the unweighted source trace,
then reconstructs the weight transport with separate prefix-event and
residue-anchor formulas. It does not call the producer's transport helpers.
Its final canonical histogram must match the certificate, including sorted
unique vectors, positive multiplicities and total mass. Conservation alone
would be insufficient; the local trace replay establishes which original
weights belong to each orbit. Tests disable producer entry points during
replay and mutate arithmetic, source and trace fields.

This is an implementation of classical weighted AHT's polynomial encoded
operation~\cite[Theorem 16]{aht-orbits}. Canonical runs prevent an implicit
expansion of the point universe. A translation fold maps old suffix breakpoints
into one period and introduces only a constant number of boundary points;
reflection also transports breakpoints directly. Static gap boundaries are
charged to the explicit pairing system. Thus the number of weight runs and
emitted histogram blocks remains polynomial in input runs, pairings and
trace length. Coordinate magnitudes can grow by summing many original points,
but their bit lengths grow only by the logarithm of the represented mass.
The dimension $d$, endpoint bits, signed weight bits and exact integer
arithmetic are all part of the cost model. Section~\ref{sec:suffix-folds}
gives the sharp native bound $R_T\le R_0+2T$, explicit mass and output
bounds, and the later optimization that restricts vector arithmetic to each
receiving interval.

The trace is needed internally even when the caller does not request a
certificate. A supplied source-bound trace can instead be checked and reused,
allowing several different weight queries on the same interval system.
This still performs exact arithmetic replay; it avoids discovering a new
orbit-reduction schedule. The cycle allowance bounds orbit discovery. Later
weight transport and independent verification are controlled by their
cooperative callbacks, not silently charged to that cycle counter.

\subsection{Three weights suffice on a torus boundary}

The validated normal arc system has one orbit per connected surface component.
Integral ownership assigns $+1$ per surface vertex, $-1$ per arc and $+1$ per
polygon, with every charge placed at an intersection point in that cell's
component. The sum is its Euler characteristic. Two more interval weights
measure intersections of its boundary with a homology basis of the ambient
boundary torus. The producer constructs that basis by a primal tree and a
disjoint dual tree. The verifier independently checks that the proposed cycles
are closed and independent modulo face boundaries using binary elimination.
Generalized triangulations with loops and repeated face edges are supported.

As proved in Section~\ref{sec:incoming-2c7f}, a component is a compressing
disc exactly when its Euler characteristic is one and its two boundary
parities are not both zero. A nonzero boundary class ensures the component
has boundary; classification then forces a disc. On a torus an essential
embedded circle has primitive slope and nonzero mod-two class. A closed
projective plane also has Euler characteristic one but zero boundary values,
so it is excluded without an irreducibility assumption.

Componentwise evaluation is necessary. Two parallel meridian discs have zero
total mod-two boundary class but contribute two discs. Conversely a
vertex-link disc and a nondisc component can yield misleading aggregate
Euler/homology values. The weighted histogram distinguishes their contributions.
The external audit includes both phenomena and one-sided components.

The public census has three modes. \texttt{disk} uses three weights;
\texttt{summary} uses five, adding polygon and boundary-point counts;
\texttt{coordinates} uses $7t$ weights and returns actual connected-component
normal coordinate vectors with multiplicities. Boundary-point count is not
the number of boundary circles. Full-coordinate mode is useful when an actual
disc vector must be supplied to a later geometric operation; the cheaper
three-weight query is sufficient when only the count is needed. All modes
can reuse the same verified unweighted trace.

\subsection{The quadrilateral core shrinks the numerical query}

For each global vertex, subtract the minimum incident triangle count times
its vertex-link vector. These links can be realized disjointly from the
remaining surface. Their boundary discs are inessential and do not affect
the compressing-disc count. Let $Y$ be the remainder and let $g$ be the gcd
of the original quadrilateral coordinates. Matching equations propagate
divisibility by $g$ from a zero triangle at each global vertex, so $P=Y/g$
is integral and admissible whenever $g>0$. If all quadrilaterals vanish,
$Y=0$ and the essential-disc count is zero without any orbit query.

The proof in Section~\ref{sec:incoming-2c7f} gives
\[
 D(x)=gD(P),\qquad
 \max_i P_i\le(4t-1)\frac{Q_{\max}}g.
\]
This identity concerns discs. It is false for the total number of components:
two copies of a one-sided normal surface can form its connected orientation
double. The count-only API consequently reports no inferred total component
count. The original unreduced census remains available for that question.

For $x=gD+(g+1)L$, where $D$ is a fixed primitive vertex-reduced meridian and
$L$ is a vertex link, the full coordinate gcd can be one. Common scaling
alone then achieves nothing. The new operation removes the $g+1$ links and
queries the fixed core $D$, returning exactly $g$ essential discs. The core
orbit endpoints have size independent of the bit length of $g$. Reading,
validating and fingerprinting $x$, computing its minima/gcd and verifying the
outer count still depend on the original input bits. The optimization removes
those bits from the expensive component query, not from the entire call.

The producer validates the original vector, constructs the core, and invokes
a source-validated three-weight census on it. The count certificate binds the
original source, all removed link multiplicities, the divisor, core vector,
core certificate and final multiplication. The verifier independently
reconstructs the minima and gcd rather than invoking the core producer.
A modified original vector is rejected even if it has the same reduced core.
This prevents replacing the original source binding with only a small-core
fingerprint.

\subsection{Native contracts and validation}

The maintained census, independent census verifier and reduced count query
are available through these entry points:
\begin{quote}\small
\path{normal_surface_components.normal_component_census}\\
\path{normal_component_verify.verify_normal_component_certificate}\\
\path{normal_disk_kernel.normal_compressing_disk_count}
\end{quote}
The last module also supplies the independent count-certificate verifier.
The generic weighted APIs are in \path{weighted_orbits} and
\path{weighted_orbit_verify}. The census can receive a supplied orbit trace;
then \texttt{max\_cycles} must be \texttt{None}, because no orbit discovery
runs. Invalid mode, cycle and periodic-rule options are rejected before
geometry work. Incomplete queries publish no component histogram, disc count
or certificate. Cancellation propagates during validation, discovery,
transport, interpretation and replay, including for false-valued callables.

The original geometry validator and orbit engine are byte-identical to baseline
\texttt{04a66388e255}. Six new runtime modules implement the extension. Normal
geometry and its exact source encoding are shared trusted reconstruction;
weight arithmetic, basis verification, component interpretation and core
arithmetic have separate checking paths. This independence is supplemented
by an external Regina oracle, which does not use producer geometry or weight
transport to obtain expected components.

The report's 30 new test methods and four integration tests pass. They include
source/summary mutations, generalized boundary bases, actual component-vector
multisets, enormous multiplicities, one-sided scaling and source-bound trace
reuse across all three modes with orbit discovery disabled. The native audit
repeats 10,000 small weighted-system comparisons across the two periodic rules,
2,500 independent checker cases, 366 mutation checks and huge endpoint tests.
Its frozen external corpus contains 1,275 vectors on 45 triangulations, with
3,825 census comparisons and replays plus 1,275 core-count comparisons and
replays. Regina explicitly expands only these bounded fixtures; the maintained
queries use compressed coordinates. The corpus includes 198 cases with discs
whose total boundary class cancels, and cases whose aggregate data suggests
a disc that is absent. The native implementations pass all these checks.

\subsection{Complete supplied-vector measurements}

The maintained suite passes all 1,084 tests with Regina in 209.265 seconds;
the 34 focused tests pass in 1.566 seconds. The native audit completes in
35.314 seconds. A separate timing run takes 280.161 seconds and completes
all 200 measured calls and 40 retained warm-ups. Ten fixtures each have two
query methods, an identical A/A copy of each, five shuffled measured rounds
and one warm-up per arm. Each timed call includes fresh input validation,
certificate production and independent source replay. Certificate serialization
and the expected-count assertion are outside the clock. Both audit and timing
freeze the same 153 source, driver, test, fixture and corpus hashes. The four
original codec, orbit and geometry dependencies match the baseline exactly.

These are comparisons between the new query interfaces. Full-coordinate
census provides more output than a three-weight disc query, and an unreduced
census provides more output than the core count-only operation. The table
therefore measures the saving from requesting sufficient information for the
disc-count task; it is not a comparison against a previously available
whole-knot recognizer.

The first family uses primitive meridian discs in layered solid tori with
$t=4,8,16,32,64,128$ tetrahedra. Every query reports one compressing disc.
The full census carries $7t$ weights and the disc census carries three.

\begin{center}\small
\begin{tabular}{@{}rrrrrr@{}}
\toprule
Tetrahedra & full ms & three ms & full/three & full A/A & three A/A \\
\midrule
4 & 4.125 & 3.065 & 1.396 & 1.002 & 1.010 \\
8 & 12.352 & 6.392 & 1.972 & 0.999 & 1.004 \\
16 & 55.524 & 17.167 & 3.179 & 1.023 & 0.974 \\
32 & 303.772 & 60.642 & 4.973 & 1.009 & 1.017 \\
64 & 2254.345 & 310.024 & 7.338 & 1.002 & 1.073 \\
128 & 16945.105 & 1698.802 & 9.975 & 1.015 & 1.001 \\
\bottomrule
\end{tabular}
\end{center}


At $t=128$, complete full-coordinate and three-weight medians are
16,945.105 and 1,698.802 ms; the median paired ratio is 9.975. Both execute
133 orbit cycles and 21,008 replay events. The saving comes from weight
construction and transport, rather than a different orbit schedule: weight
dimension falls from 896 to three and maximum constant-weight runs from
385 to eight. The compact serialized certificate falls from 2,504,362 to
1,777,349 bytes; the large common orbit trace remains. Across this family,
full A/A paired medians range from 0.999 to 1.023 and three-weight A/A from
0.974 to 1.073. These controls are much smaller than the larger dimension
savings. No asymptotic exponent is fitted to these six timing points.

The second family fixes $t=16$ and supplies
$x=2^bD+(2^b+1)L$ for $b=64,512,4096,32768$. Its full coordinate gcd is one,
while every query must count exactly $2^b$ essential discs. Both arms already
use three weights. The comparison isolates vertex-link removal and
quadrilateral-content division, including the original-source checks and
outer count proof.

\begin{center}\small
\begin{tabular}{@{}rrrrrr@{}}
\toprule
$b$, $g=2^b$ & raw ms & core ms & raw/core & raw A/A & core A/A \\
\midrule
64 & 42.837 & 21.948 & 1.952 & 1.021 & 0.996 \\
512 & 83.938 & 41.735 & 2.108 & 0.999 & 0.995 \\
4096 & 116.878 & 43.794 & 2.688 & 1.018 & 0.981 \\
32768 & 296.942 & 35.789 & 5.620 & 1.041 & 0.967 \\
\bottomrule
\end{tabular}
\end{center}


For every scale the core has 11-bit coordinates and 14-bit orbit endpoints.
The unreduced query uses 39 cycles, 819 replay events and at most 51 weight
runs; the core uses 21 cycles, 400 events and at most eight runs. At
$b=32768$, the original maximum coordinate length is 32,779 bits and its
unreduced orbit endpoints have 32,782 bits. The compact proof shrinks from
4,523,182 to 56,757 bytes, while still binding the original source. The
reduced orbit work is independent of $b$ for this family; the outer reading,
validation, arithmetic and fingerprint work remains sensitive to $b$.

The largest core case has appreciable timing variability. Its five raw/core
ratios are 5.187, 4.784, 8.899, 5.620 and 5.977; their median is 5.620.
The separate arm medians, 296.942 and 35.789 ms, must not be divided and
reported as that paired ratio. The identical core control has a 60.289 ms
median, and individual raw/core timings change substantially across rounds,
even though their median A/A ratios are 1.041 and 0.967. Across the core
family the raw A/A medians range from 0.999 to 1.041 and core A/A from
0.967 to 0.996. Median controls alone do not capture the observed variation.
The raw samples and warm-ups are retained; deterministic operation counts,
bit sizes and certificate sizes substantiate the reduction independently of
the precise elapsed-time factor.

These measurements neither include finding a normal vector nor establish a
uniform saving on arbitrary triangulations. Small already primitive vectors
can incur reduction overhead. No automatic knot-decision stage is added by
this integration, and a zero count is not evidence that other disc vectors
do not exist.


Reproduction uses \path{../fast/weighted_research/native.py}, with
\texttt{audit} and \texttt{benchmark} modes. It pins the unchanged native
dependencies and all runtime, driver, test, fixture and corpus hashes before
and after each run. Data are in \path{data/weighted-components-*}; the delivered
report is unchanged. The two-weight ambient-homology refinement remains
unimplemented. Section~\ref{sec:diagram-exterior} now provides certified
diagram-to-exterior construction for a canonical unsimplified triangulation.
Useful normal-vector discovery and hierarchy moves with controlled encoded
cost remain global tasks. Efficiently counting discs in one supplied vector
does not settle those tasks or establish general quasipolynomial recognition.
